\appendix

\section{Formal architecture and notation}\label{app:model}

We restate the architecture only to fix implementation details. At layer $\ell$, one causal
attention head maps representations $H_0,\ldots,H_p$ to
\[
 q_p=W_QH_p+b_Q,\qquad k_s=W_KH_s+b_K,\qquad v_s=W_VH_s+b_V,
\]
then returns the average of $v_s$ over score maximizers of $\langle q_p,k_s\rangle$ for
$s\leq p$. A multi-head layer sums head outputs, which may occupy disjoint coordinate slots. A
position-wise constant-depth ReLU network follows attention. Residual connections may preserve
input slots. Token embeddings are trainable. Positional embeddings are fixed. A decoding matrix
maps the last representation to vocabulary logits. This is the model of
\citet[Section~2.1]{yang2026tight}.

The architecture is fixed for each $(g,r)$. Fixing a trainable coordinate to a stated constant means
restricting the full architecture class to a subclass. This is valid for a VC lower bound. All widths
between Transformer blocks are bounded by a universal constant. The optional group lookup uses a
width-$O(g)$ internal FFN layer, which is permitted while keeping the block representation
dimension constant.

\subsection{Detailed end-to-end upper bound}\label{app:upper}

Let $\Pi_A(n)$ be the largest number of joint next-token patterns that $\mathcal F_A$ realizes on
any $n$ fixed contexts of length at most $T_{\rm ctx}$. Equation~\eqref{eq:yang-growth} states
\[
 \log\Pi_A(n)\leq C_0WL\log(C_1nT_{\rm ctx}WLK).
 \tag{A.1}
\]
Fix $n$ prompts. Partition parameter space by the $n$ tokens emitted on call one. There are at most
$\Pi_A(n)$ cells. Inside any cell, all prefixes presented on call two are fixed. Intersect that cell
with the global next-token pattern partition for those prefixes. This creates at most
$\Pi_A(n)$ subcells. Continuing inductively yields at most $\Pi_A(n)^M$ full-trajectory cells.
The final-token pattern is fixed on each cell. Hence shattering implies
\[
 2^n\leq\Pi_A(n)^M,
 \qquad
 n\leq C_0MWL\log(C_1nT_{\rm ctx}WLK).
 \tag{A.2}
\]
If $n\leq a\log(cn)$ with $a,c\geq2$, then $n\leq C'a\log(C'ac)$. Applying this with
$a=C_0MWL$ gives Proposition~\ref{prop:upper}.

This recursion is needed because the prefixes at later calls are not fixed globally. It is the same
reason that applying the one-pass VC bound to only the original prompts would be invalid.

\section{Exact Transformer primitives}\label{app:primitives}

The construction uses four elementary identities.

\begin{lemma}[Binary ReLU gate]\label{lem:gate}
If $b\in\{0,1\}$ and $x\in[0,C]$, then
\[
 bx=\operatorname{ReLU}(x+C(b-1)).
 \tag{B.1}
\]
For $x\in[0,1]$, one may take $C=1$.
\end{lemma}

\begin{proof}
For $b=1$ the argument is $x$. For $b=0$ the argument is nonpositive.
\end{proof}

Thus a token-type bit can gate any bounded fixed positional feature. All bounds are finite because
$p\leq T_{\rm ctx}-1$. Negative features are represented by positive and negative parts or by an
affine combination of gated nonnegative features.

\begin{lemma}[Typed average]\label{lem:average}
Suppose visible positions carry a type flag $u_s\in\{0,1\}$ and value $v_s$, and at least one has
$u_s=1$. A hard-attention head with constant query one and key $u_s$ returns the average of $v_s$
over positions with $u_s=1$.
\end{lemma}

\begin{lemma}[Quadratic address]\label{lem:address}
For an integer target $q$ and integer positions $p$, the score
\[
 2qp-p^2=q^2-(p-q)^2
 \tag{B.2}
\]
has the unique maximizer $p=q$ whenever that position is visible.
\end{lemma}

Both lemmas follow directly from the definition of attention. The address is an inner product of
query $(2q,1)$ and key $(p,-p^2)$.

\begin{lemma}[Binary selector]\label{lem:selector}
For $u,v,F\in\{0,1\}$,
\[
 \operatorname{ReLU}(u-F)+\operatorname{ReLU}(v+F-1)
 =\begin{cases}u,&F=0,\\v,&F=1.\end{cases}
 \tag{B.3}
\]
\end{lemma}

\section{Full tape construction}\label{app:construction}

Fix $g,r\geq2$ and let $\ell=\lceil\log_2(r+1)\rceil$. The main alphabet is
\[
 \{G_1,\ldots,G_g,I_0,I_1,A,S,B_0,B_1,P\}.
\]
Token $P$ is used only for optional padding. Prompt $(j,i)$ is
\[
 G_j,\ I_{i_0},\ldots,I_{i_{\ell-1}},\ A,\ S,
 \tag{C.1}
\]
where $i=\sum_{q=0}^{\ell-1}i_q2^q$. All prompts have length
$T=\ell+3$. Positions are zero indexed. Thus the first generated bit occupies position $T$ and
bit $s$ occupies position $T+s-1$.

We allocate disjoint constant-dimensional coordinate slots. A residual path preserves the slots that
later blocks need. Since head outputs occupy disjoint slots, describing them separately is
equivalent to giving the corresponding sparse projection matrices.

\subsection{Embeddings and fixed positional features}

Token embeddings contain the following flags and values:
\begin{center}
\begin{tabular}{lcccccc}
\toprule
token & group & index & bit type & token bit & fallback & tape value\\
\midrule
$G_j$ & 1 & 0 & 0 & 0 & 1 & $\theta_j$\\
$I_0$ & 0 & 1 & 0 & 0 & 1 & 0\\
$I_1$ & 0 & 1 & 0 & 1 & 1 & 0\\
$A$   & 0 & 0 & 0 & 0 & 1 & 0\\
$S$   & 0 & 0 & 1 & 0 & 0 & 0\\
$B_0$ & 0 & 0 & 1 & 0 & 0 & 0\\
$B_1$ & 0 & 0 & 1 & 1 & 0 & 0\\
$P$   & 0 & 0 & 0 & 0 & 1 & 0\\
\bottomrule
\end{tabular}
\end{center}
The fallback column is not essential. Every non-bit token has zero comparator key and zero
comparator value, which is enough.

The fixed positional embedding supplies $1,p,p^2$ at every position. At index position $q+1$ it
also supplies $2^q$. At positions $p\geq T-1$, it supplies
\[
 s(p)=p-T+2,
 \qquad d(p)=2^{-s(p)},
 \tag{C.2}
\]
and at generated-bit position $p=T+s-1$ it supplies $2^{-s}$. Finally it supplies
\[
 F(p)=\1\{p=T+r-1\},
 \tag{C.3}
\]
which is one exactly at the position that is current when the final answer is generated. Values of
these auxiliary coordinates outside their intended position ranges may be zero.

The label assignment chooses
\[
 \theta_j=\sum_{s=1}^rY_{j,s}2^{-s}+2^{-(r+1)}.
 \tag{C.4}
\]
Every other embedding and weight is fixed.

\subsection{Block one: typed features}\label{app:block1}

Set the attention sublayer to the identity through a residual. Its position-wise FFN uses
Lemma~\ref{lem:gate} to form:
\begin{align*}
 &u^{\rm bit},\quad u^{\rm group},\quad u^{\rm index},\\
 &u^{\rm bit}p,\quad u^{\rm bit}p^2,\quad
 u^{\rm bit}s(p),\quad u^{\rm bit}2^{-s(p)},\\
 &v^{\rm prefix}_p
 =\begin{cases}Y_{j,s}2^{-s},&p=T+s-1,\\0,&\text{otherwise},\end{cases}\\
 &v^{\rm index}_p
 =\begin{cases}i_q2^q,&p=q+1,\\0,&\text{otherwise}.
 \end{cases}
\end{align*}
For $v^{\rm prefix}$, equation~\eqref{eq:relu-product} uses the token bit and the fixed dyadic
feature. For $v^{\rm index}$, use Lemma~\ref{lem:gate} with $C\geq2^{\ell-1}$.

All features are exact on the finite token-position domain. A constant number of ReLU units per
feature suffices. The width and number of trainable entries are independent of $g$ and $r$.

\subsection{Block two: group, prompt index, and emitted prefix}\label{app:block2}

Three heads apply Lemma~\ref{lem:average}.
\begin{enumerate}[leftmargin=*]
 \item The group head scores $u^{\rm group}$ and returns the unique value $\theta_j$.
 \item The index head scores $u^{\rm index}$ and returns
 $\ell^{-1}\sum_qi_q2^q=i/\ell$. A fixed affine map in the FFN multiplies by $\ell$.
 \item The prefix head scores $u^{\rm bit}$. While tape bit $t$ is about to be emitted, the
 selected positions are $S$ and the preceding $t-1$ bit tokens. Their count is $t$, so its output is
 \[
 a_t=\frac1t\sum_{s<t}Y_{j,s}2^{-s}=\frac{P_{t-1}}t.
 \tag{C.5}
 \]
\end{enumerate}
The FFN preserves $\theta_j,a_t,i,p_*,p_*^2$ in separate slots at the current position. It does not
form $t a_t$.

\subsection{Block three: next bit and requested bit}\label{app:block3}

At the current position $p_*$, define comparison query
\[
 q^{\rm cmp}_{p_*}=(\theta_j,a_t,1,p_*,p_*^2).
 \tag{C.6}
\]
At position $p$, define comparison key
\[
 k^{\rm cmp}_p=u^{\rm bit}_p
 \bigl(1,-s(p),-2^{-s(p)}-2p^2,4p,-2\bigr).
 \tag{C.7}
\]
The dot product is equation~\eqref{eq:comparison-score}. Non-bit positions have key zero and score
zero. Their comparison value is zero. Bit-type positions have comparison value one.

At tape step $t$, the current bit-type position has $s(p_*)=t$, so by equation~\eqref{eq:tape-decode}
its score has sign $Y_{j,t}$. For an earlier bit-type position $p<p_*$,
\[
 S(p)\leq\theta_j-2(p-p_*)^2<1-2<0,
 \tag{C.8}
\]
where we dropped two nonpositive terms. If $Y_{j,t}=1$, the current position is the unique positive
maximizer and the head returns one. If $Y_{j,t}=0$, every bit-type score is negative and every
non-bit position maximizes at zero with value zero. The head returns zero. Call this output $c_t$.

The retrieval head has query
\[
 q^{\rm ret}_{p_*}=(2(T+i-1),1)
\]
and key $k^{\rm ret}_p=(p,-p^2)$. Its value is the token bit on $B_0,B_1$ and zero on prompt
tokens. At the final call, all tape positions are visible. Lemma~\ref{lem:address} gives output
$d=Y_{j,i}$.

The block-three FFN applies Lemma~\ref{lem:selector} with $F(p_*)$:
\[
 z=\operatorname{ReLU}(c_t-F(p_*))
 +\operatorname{ReLU}(d+F(p_*)-1).
 \tag{C.9}
\]
It equals $c_t$ for the first $r$ calls and $d$ on the final call. Decoder logits
$\operatorname{logit}(B_1)=z-1/2$ and
$\operatorname{logit}(B_0)=1/2-z$ emit the desired bit. All other logits are fixed below both.

Equations~(C.4)--(C.9) prove by induction that the generated trace is
\[
 Y_{j,1},\ldots,Y_{j,r},Y_{j,i}.
\]

\subsection{Parameter count}\label{app:count}

The hidden dimension, number of heads, and all backbone matrix shapes are constant. The group-token
embedding table and vocabulary decoder contain $O(g)$ entries. Every other table or matrix has
$O(g)$ entries only because the vocabulary has $g+O(1)$ rows. Therefore $W(A_{g,r})\leq Cg$.
The construction uses three attention-FFN blocks and $T=\ell+3=O(\log r)$ prompt positions.

\section{Constant-alphabet compilation}\label{app:constant-alphabet}

Replace $G_j$ by $m=\lceil\log_2g\rceil$ binary prompt tokens. Fixed positional weights and the
index aggregation of Appendix~\ref{app:block2} recover integer $j\in\{0,\ldots,g-1\}$. Let the
desired tape scalars be $\theta_0,\ldots,\theta_{g-1}$. Define ReLU hat functions on this integer
domain by
\begin{align}
 \phi_0(x)&=\operatorname{ReLU}(1-x),\nonumber\\
 \phi_q(x)&=\operatorname{ReLU}(x-q+1)-2\operatorname{ReLU}(x-q)
 +\operatorname{ReLU}(x-q-1), &&1\leq q\leq g-2,\nonumber\\
 \phi_{g-1}(x)&=\operatorname{ReLU}(x-g+2)-\operatorname{ReLU}(x-g+1).
 \tag{D.1}
\end{align}
For integer $j,q\in\{0,\ldots,g-1\}$, $\phi_q(j)=\1\{j=q\}$. Therefore
\[
 h(j)=\sum_{q=0}^{g-1}\theta_q\phi_q(j)=\theta_j.
 \tag{D.2}
\]
A width-$O(g)$ ReLU FFN computes these hats and writes one scalar output. Crucially, each tape
scalar $\theta_q$ is one direct output-layer parameter rather than a difference of several tape
parameters. The network has $O(g)$ parameters and constant input and output dimension.

The alphabet can now be
\[
 \{0,1,A,S,B_0,B_1,P\}.
\]
Position ranges distinguish the group bits from the requested-index bits. The prompt length becomes
$O(\log g+\log r)$, and the rest of the construction is unchanged.

\section{Finite trainable precision}\label{app:precision}

For completeness, a $b$-bit coordinate codebook may be different for each parameter. The upper
bound uses only
\[
 |Q_1\times\cdots\times Q_W|\leq2^{bW}.
 \tag{E.1}
\]
For the lower bound, choose $r\leq b-1$ and use the codebook
\[
 Q=\left\{\sum_{s=1}^ry_s2^{-s}+2^{-(r+1)}:y\in\{0,1\}^r\right\}.
 \tag{E.2}
\]
It has $2^r\leq2^b$ entries. Each entry has a terminating binary expansion with $r+1$ fractional
bits. In the constant-alphabet compiler, these values occur directly as the output weights in
equation~(D.2). Every varying coordinate therefore uses this same $2^r$-element codebook.

If the prescribed rollout length is $M>r+1$, add token $P$. A fixed positional flag emits $P$ on
calls $r+1,\ldots,M-1$ and selects the retrieval head on call $M$. The tape prefix remains visible,
so the final answer is unchanged. This establishes Theorem~\ref{thm:precision} for every $M$.

The definition counts selectable trainable values. If every fixed coefficient must also use a
common signed fixed-point format, the largest values are the index weights $2^q\leq r+1$ and the
softmax scale from Appendix~\ref{app:softmax}. Their binary descriptions use
$O(r+\log T_{\rm ctx})$ bits.

\section{Full proof of the supervision gap}\label{app:gap}

The end-to-end map of $\mathcal T_{g,r}$ is exactly
\[
 (j,i)\mapsto Y_{j,i},
 \qquad Y\in\{0,1\}^{g\times r}.
 \tag{F.1}
\]
It is the class of all binary functions on $gr$ points and has VC dimension $gr$. The optimal
realizable fixed-confidence PAC rate is $\Theta(gr/\epsilon)$
\citep{hanneke2016optimal}.

For the trace upper bound, let $p_j=\sum_iD(j,i)$. After $n$ samples, memorize $Y_j$ for every
observed group. On an unseen group, use any default tape. The final-token error is at most
\[
 U_n=\sum_{j=1}^gp_j\1\{N_j=0\}.
\]
For $x\in[0,1]$, $x(1-x)^n\leq1/(e(n+1))$. Therefore
\[
 \E U_n=\sum_jp_j(1-p_j)^n\leq\frac{g}{e(n+1)}.
 \tag{F.2}
\]
At confidence $3/4$, Markov's inequality makes $U_n\leq\epsilon$ once
$n\geq4g/(e\epsilon)$. Any other fixed confidence changes only the constant.

For the lower bound, restrict equation~(F.1) to tapes
$Y_{j,1}=\cdots=Y_{j,r}=z_j$. Both the full trace and final answer reveal only arbitrary group bit
$z_j$. The standard rare-points construction for a $g$-point binary domain requires
$\Omega(g/\epsilon)$ examples at constant confidence. This proves both sides of
Theorem~\ref{thm:gap}.

The usual $gr$-point rare-points distribution can be used for the end-to-end lower bound. Since
equation~(F.2) holds for every distribution, the trace upper bound holds on that same distribution.
Thus the factor-$\Theta(r)$ comparison does not rely on changing the prompt law between regimes.

The learner is allowed to be improper in both regimes. The separation therefore does not arise
from a restriction on the training algorithm.

\section{Softmax compilation}\label{app:softmax}

We first give a general finite-domain observation. Suppose a hard-attention head has a maximizer
set $M(x)$ on input $x$, and every nonmaximizer is below the maximum by at least $\gamma>0$.
If all value norms are at most $V$ and at most $N$ positions are visible, softmax with inverse
temperature $\beta$ differs from the uniform average on $M(x)$ by at most
\[
 2VN e^{-\beta\gamma}.
 \tag{G.1}
\]
This follows by bounding the total nonmaximizer mass relative to one maximizing exponential. It
also applies when the maximizers have different values, as in the prefix-average head.

In the tape construction, group, type, and exact-address heads have constant gaps. The comparison
head has gap at least
\[
 \gamma_{\rm cmp}=2^{-(r+1)}
 \tag{G.2}
\]
between the current score and zero in the positive case, or between zero and every bit score in the
negative case. Earlier bit positions have an additional gap greater than one. The selector heads
must approximate $\theta_j$ and $a_t$ accurately enough that the perturbed comparison residual
retains margin at least $\gamma_{\rm cmp}/2$.

All fixed affine and ReLU maps in three blocks have a finite Lipschitz constant
$\operatorname{poly}(r,T_{\rm ctx})$ on the relevant bounded states. Choose selector inverse
temperatures so equation~(G.1) is at most
$\gamma_{\rm cmp}/\operatorname{poly}(r,T_{\rm ctx})$. It suffices to take
\[
 \beta_{\rm sel}=O(r+\log T_{\rm ctx}).
 \tag{G.3}
\]
After this perturbation, take
\[
 \beta_{\rm cmp}=2^{O(r)}\operatorname{polylog}(T_{\rm ctx})
 \tag{G.4}
\]
so the comparison output is on the correct side of $1/2$. Constant-gap retrieval needs only
$O(\log T_{\rm ctx})$ inverse temperature.

Induct on autoregressive calls. If the preceding softmax outputs decode to the correct discrete
tokens, the current context is exactly the hard construction's context. Equations~(G.1)--(G.4)
then make the next logit choose the correct token. The base call is identical. Thus all $r+1$
tokens are preserved for every prompt and every one of the finitely many tape assignments. The
largest scale in equation~(G.4) has $O(r+\log T_{\rm ctx})$ bits in binary scientific notation.

\section{Executable audits}\label{app:audit}

The supplement contains two dependency-free Python programs.

\paragraph{Exact rational audit.}
\texttt{audit\_tape.py} uses \texttt{fractions.Fraction}. For every tape through length 14, it
checks:
\begin{enumerate}[leftmargin=*]
 \item sequential recovery by equation~\eqref{eq:tape-decode}
 \item exact minimum margin $2^{-(r+1)}$
 \item the full comparison score ordering in equation~\eqref{eq:comparison-score}
 \item unique quadratic retrieval for every requested coordinate
 \item the constant-alphabet hat lookup in equations~(D.1) and (D.2).
\end{enumerate}
This enumerates $\sum_{r=1}^{14}2^r=32766$ tapes.

\paragraph{Softmax audit.}
\texttt{audit\_softmax.py} replaces each hard selector, prefix average, comparison, and retrieval
with stable finite-temperature softmax. It uses explicit scales following equations~(G.3) and
(G.4), then enumerates every tape and requested coordinate through length 12. All 8190 tapes and
their rollouts pass.

The programs print one line per tape length and terminate with assertion failures if any identity,
score ordering, emitted token, or retrieved coordinate is wrong. They are audits of the scalar
construction rather than a training experiment.
